% Part I > What formalization fixes: the problems the community response
% exposed, what a machine-checked certificate (used correctly) does about each,
% where it stops (Navier--Stokes), and the open problems handed to Part II.
% Facts and sources: research/01 (s5, Navier--Stokes and the OpenAI audit),
% 02 (Palomar, Comparator, Prove2Me), 03 (s5, what formalization can and
% can't do).
\section{What formalization fixes}

% [2 min] Introduces AI-assisted formalization as a concept: the pipeline
% (human statement and blueprint -> AI formalizes and proves -> kernel checks
% -> registered certificate), then what it addresses from the community-response
% section and what it leaves open. Replaces the earlier problem/certificate
% table (same content, Kyle asked for a non-table slide).
\begin{frame}{AI-assisted formalization}
  \centering
  \begin{tikzpicture}[x=1cm, y=1cm, font=\sffamily\footnotesize]
    \useasboundingbox (0,0.1) rectangle (14.3,-5.9);
    % ---- definition --------------------------------------------------------
    \node[anchor=north west, text width=14.2cm, align=flush left, inner sep=0pt, font=\sffamily\small]
      at (0,0.05)
      {AI translates mathematics into a proof assistant such as Lean and fills in the proofs.
       A small trusted kernel checks every step. Humans write, or review, what is being proved.};
    % ---- pipeline ----------------------------------------------------------
    \node[stage=PennBlue, minimum width=2.9cm, minimum height=1.05cm] (p1) at (1.65,-1.5)
      {Statement \& blueprint\\[-1pt]{\scriptsize\color{black!60}human}};
    \node[stage=tlPurple, minimum width=2.9cm, minimum height=1.05cm] (p2) at (5.35,-1.5)
      {Formalize \& prove\\[-1pt]{\scriptsize\color{black!60}AI agents}};
    \node[stage=PennGreen, minimum width=2.9cm, minimum height=1.05cm] (p3) at (9.05,-1.5)
      {Kernel check\\[-1pt]{\scriptsize\color{black!60}Lean}};
    \node[stage=PennGold, minimum width=2.9cm, minimum height=1.05cm] (p4) at (12.75,-1.5)
      {Certificate\\[-1pt]{\scriptsize\color{black!60}registry, e.g.\ Palomar}};
    \foreach \u/\v in {p1/p2, p2/p3, p3/p4} \draw[flow] (\u) -- (\v);
    \draw[flow=PennBlue, dashed] (p3.south) .. controls +(0,-0.55) and +(0,-0.55) .. (p1.south);
    \node[note, fill=white, inner sep=1pt] at (5.35,-2.48) {a human checks that the formal statement matches the paper};
    % ---- what it addresses / what stays open --------------------------------
    \node[anchor=north west, inner sep=0pt, font=\sffamily\small\bfseries, text=PennGreen] at (0,-2.95)
      {What it addresses};
    \node[anchor=north west, inner sep=0pt, font=\sffamily\small\bfseries, text=PennRed] at (7.35,-2.95)
      {What stays open};
    \node[anchor=north west, text width=6.85cm, align=flush left, inner sep=0pt, font=\sffamily\scriptsize] at (0,-3.4) {%
      \begin{itemize}\setlength{\itemsep}{1pt}
        \item \textbf{Review capacity.} Referees review a page of statements, not 60 pages of proof.
        \item \textbf{Undisclosed AI.} The check does not depend on who, or what, wrote the proof.
        \item \textbf{Bypassed review.} Verify and register a result before announcing it.
        \item \textbf{Credit.} Registries record what produced each result.
      \end{itemize}};
    \node[anchor=north west, text width=6.85cm, align=flush left, inner sep=0pt, font=\sffamily\scriptsize] at (7.35,-3.4) {%
      \begin{itemize}\setlength{\itemsep}{1pt}
        \item \textbf{Statement fidelity.} Compiling is not saying the same thing (Navier--Stokes).
        \item \textbf{Understanding.} Millions of machine-written lines explain nothing.
        \item \textbf{Significance and ideas.} Neither is a formal property.
        \item \textbf{Access and coverage.} Frontier compute; thin libraries in PDE and applied probability.
      \end{itemize}};
  \end{tikzpicture}
  \slidesources{Tao, \emph{Palomar}, blog, 18 Aug 2026 \textbullet{} Bastounis, Circelli \& Hansen, arXiv 2610.08144 \textbullet{} research notes 01--03}
\end{frame}

% [1.5 min] Navier--Stokes revisited: the certificate compiled, and the code
% still says something other than the paper. They do NOT claim the paper is
% wrong; the argument was about credit, not correctness.
\begin{frame}{The paper and the code disagree}
  \small\sffamily
  \begin{columns}[T]
    \begin{column}{0.54\textwidth}
      \textbf{Navier--Stokes} (Bastounis, Circelli \& Hansen, Oct 2026)
      \begin{itemize}\footnotesize
        \item The paper's estimate (8.19) needs $m{+}4$ derivatives; the matching Lean theorem assumes $m{+}5$, and the pattern repeats.
        \item The paper's pressure bound depends on one quantity; the Lean bound depends on two, and is proved by a different route.
        \item They do not claim the paper is wrong.
      \end{itemize}
      \medskip
      \textbf{OpenAI's 719 manuscripts} (community audit of 245 formalized families)
      \begin{itemize}\footnotesize
        \item About 210 consistent at the level of statements; no theorem shown false.
        \item 21 headline results marked ``(Lean)'' were not formalized; 5 formalized only a supporting lemma.
      \end{itemize}
    \end{column}
    \begin{column}{0.42\textwidth}
      \vspace{0.4cm}
      \begin{tikzpicture}
        \node[draw=PennRed!45, fill=PennRed!3, line width=0.6pt, rounded corners=2pt,
              text width=5.3cm, inner sep=0.3cm, align=left, font=\rmfamily\itshape\small] (q)
          {Lean acceptance (compilation) of translation does not imply semantic preservation.};
        \node[anchor=north west, font=\sffamily\scriptsize, text=black!60, inner sep=0pt]
          at ([yshift=-0.15cm]q.south west) {--- Bastounis, Circelli \& Hansen, arXiv 2610.08144};
      \end{tikzpicture}
      \vspace{0.5cm}

      {\footnotesize A certificate proves \emph{a} theorem. Whether it is \emph{the} theorem is a human question.}
    \end{column}
  \end{columns}
  \slidesources{Bastounis, Circelli \& Hansen, arXiv 2610.08144 \textbullet{} borism/openai-math-slop, issue \#2 \textbullet{} research note 01, \S5}
\end{frame}

% [2 min] The three layers: posed here, answered in Part II.
\begin{frame}{Three layers of checking}
  \centering
  \begin{tikzpicture}[x=1cm, y=1cm]
    \useasboundingbox (0,0.1) rectangle (14.2,-6.0);
    \node[layer=PennGreen, anchor=north west] (k) at (3.6,0)
      {\textbf{Does Lean accept the proof?} Cheap and mechanical. Watch for \texttt{sorry}, extra axioms and kernel bugs.\hfill{\stdone}};
    \node[lbl=PennGreen] at (3.35,-0.72) {1. Kernel};
    \node[layer=PennGold, anchor=north west] (s) at (3.6,-1.8)
      {\textbf{Does the formal statement say what the paper claims?} Scarce: ``will never be solved purely in software'' (Kontorovich). Navier--Stokes drifted here.\hfill{\stopen}};
    \node[lbl=PennGold!85!black] at (3.35,-2.52) {2. Statement};
    \node[layer=PennRed, anchor=north west] (u) at (3.6,-3.6)
      {\textbf{Does anyone understand the proof, and does it match the written one?} Scarce. Machine proofs run to millions of lines.\hfill{\stopen}};
    \node[lbl=PennRed] at (3.35,-4.32) {3. Understanding};
    \node[anchor=north, font=\sffamily\small, text width=12cm, align=center] at (7.8,-5.4)
      {Layer 1 is solved. Layers 2 and 3 are the open problems for AI-assisted formalization: Part II.};
  \end{tikzpicture}
  \slidesources{Research note 01, \S0 \textbullet{} Kontorovich, via AcadMathSci, Oct 2026 (\S2.2) \textbullet{} Bastounis, Circelli \& Hansen, arXiv 2610.08144}
\end{frame}
